\section{Old\_intro}\label{app:old-intro}

\begin{example}[A random bit generator]\label{ex:rbg}
The program $P$ generates a random bit: it resets a qubit $q$, puts it in
superposition with the Hadamard gate $H$, measures it in the computational
basis and emits the outcome on the channel $c$. It then waits for a request
on the channel $r$ before generating another bit. In the examples we omit the
tags of the prefixes.
\[
P \ =\ \rec{X}{\ \mathsf{init}(q).\,H(q).\,\M(q \rhd x).\,
       (c!x \parallel r?y.X)}
\]
The test $T$ reads the bit. If the bit is good, here if it is $0$, the test
succeeds, which is written $\mathbf{1}$, and stops; otherwise it asks $P$ to
start again, and waits for the next bit.
\[
T \ =\ \rec{Y}{\ c?z.\,\ite{z = 0}{\mathbf{1}}{(r!0 \parallel Y)}}
\]
The system is $S = P \parallel T$, run on a memory
$\rho$ that contains the qubit $q$.

One round of $S$ goes as follows. The two recursions are unfolded, in any
order. The program then reduces alone three times, for the initialisation,
the gate and the measurement; these reductions change the memory, and the
measurement has two reductions, one for each outcome $m$, each keeping half
of the trace of the memory. The bit is then communicated from $P$ to $T$. If
it is $0$, the test becomes $\mathbf{1}$ and the system stops, with a
successful test; otherwise the request is communicated from $T$ to $P$ and a
new round starts.
\qed
\end{example}

\Cref{ex:rbg} uses every feature of the calculus that matters for the Harmony
lemma. The output $c!x$ is asynchronous: it has no continuation, and is put
in parallel with the rest of the program. The initialisation is what allows
the same qubit to be used in every round. The recursion makes the system run
for an unbounded number of rounds, although it stops with probability $1$.
And the tags, which we left out, give a name to every reduction above: a
reduction is indexed by the tag of the prefix that moves, or by the pair of
the tags of the output and of the input for a communication, so that a
sequence of such choices describes one execution of the system up to the
outcomes of the measurements.

\begin{example}[Grover search with a verifier]\label{ex:grover}
The generator of \cref{ex:rbg} is the simplest instance of a common pattern, a quantum
computation whose result is checked classically and which is run again when
the check fails. Grover search is another instance. Let $w \in \{0,1,2,3\}$
be the marked element of a search space of size $4$, let $O_w$ be the oracle
that flips the phase of $|w\rangle$, and let $D$ be the diffusion gate, both
on two qubits $\tilde q = q_1, q_2$. The program $G$ prepares the uniform
superposition, applies one Grover iteration, measures the two qubits and
emits the outcome, a number between $0$ and $3$:
\[
G \ =\ \rec{X}{\ \mathsf{init}(q_1).\,\mathsf{init}(q_2).\,H^{\otimes 2}(\tilde q).\,
       O_w(\tilde q).\,D(\tilde q).\,\M(\tilde q \rhd x).\,
       (c!x \parallel r?y.X)}
\]
The verifier $V$ knows how to recognise the marked element, and asks for
another run when the answer is wrong:
\[
V \ =\ \rec{Y}{\ c?z.\,\ite{z = w}{\mathbf{1}}{(r!0 \parallel Y)}}
\]
The system $G \parallel V$ has the same reductions as $P \parallel T$, with
two more quantum steps in each round. On two qubits, one iteration is enough:
the measurement still has a reduction for each of its four outcomes, but the
memory reached for the outcomes $m \neq w$ has trace zero, so that the
verifier succeeds in the first round with probability $1$. For $n$ qubits,
the oracle and the diffusion are repeated about $\frac{\pi}{4}\sqrt{2^n}$
times; this number is fixed in the term, the recursion of \lqCCS having no
parameter and the expressions no arithmetic. A wrong answer then has a
non-zero probability, and the request for another run is really used.
\qed
\end{example}
